\documentclass[12pt]{article}

\usepackage{sbc-template}
\usepackage{graphicx,url}
\usepackage[utf8]{inputenc}
\usepackage[brazil]{babel}
\usepackage{float} % PACOTE ADICIONADO PARA FIXAR AS IMAGENS COM [H]

\sloppy

\title{Trabalho da Disciplina Gerencia de Configuração de Software}

\author{Kimberlly L. Stachelski Canhete\inst{1}, Rafaela Santos S. Orrico\inst{2}, Aluno(a)\inst{3}, Aluno(a)\inst{4}, Aluno(a)\inst{5}}


\address{Faculdade de Computação -- Universidade Federal de Mato Grosso do Sul (UFMS)\\
  79070-900 -- Campo Grande -- MS -- Brasil
  \email{\{kimberlly.stachelski, rafaela.orrico, passaporte.ufms\}@ufms.br}
}

\begin{document}

\maketitle

\begin{abstract}
    This article addresses the Software Configuration Management project, utilizing Git as its primary tool.
\end{abstract}

\begin{resumo}
  Este artigo aborda o trabalho de Gerência de Configuração de Software, utilizando o Git como ferramenta principal
\end{resumo}

\section{Introdução}
Este artigo mostrará como foi feito o trabalho da disciplina Gerencia de Configuração de Software, a ideia é mostrar o passo a passo e o que cada integrante do grupo contribuiu.

\section{Integrantes}
Esta seção mostrará se os integrantes entregaram o que foi solicitado pela docente da disciplina, assim como se foi feito tudo ou parcialmente e quais partes foi concluida com sucesso, a ideia geral desta seção é mostrar como cada integrante participou da realização do trabalho.

\subsection{Equipe}
\begin{enumerate}
    \item {Modelo Prioridade ou Modelo Severidade}
    \\\texttt{R: A escolha da equipe se deu pelo modelo de prioriodade devido a maior facilidade da classificação e o mesmo ser de facil entedimento.}
    
    \item {Modelo de Ramificação}
    \\ \texttt{R:}
    
    \item {Padrão Semantico}
    \\ \texttt{R:}
\end{enumerate}

\subsection{Kimberlly L. Stachelski Canhete}
\subsubsection{Gerenciamento de Mudanças}
    \begin{enumerate}
        \item {[x] Criar seis issues;}
        \item {[x] Reagir e comentar uma issue cadastrada por outra pessoa da equipe. O comentário deve conter informação útil para a possível solução da issue;}
        \item {[x] Revisar uma issue em status "closed"\ de outra pessoa;}
        \item {[x] Registrar como comentário na issue revisada, os ajustes que precisam ser feitos.}
    \end{enumerate}

\subsubsection{Controle de Versões}
    \begin{enumerate}
        \item {[x] Criar um arquivo dentro do projeto;} 
        \item {[x] Alterar um arquivo dentro do projeto;}
        \item {[x] Revisar e aprovar um pull-request;}
        \item {[x] Realizar uma alteração que resolva duas issues cadastradas por outro membro;}
        \item {[ ] Resolver, pelo menos, dois conflitos de merge.}
    \end{enumerate}

\subsection{Murilo Barbiere}
    \subsubsection{Gerenciamento de Mudanças}
        \begin{enumerate}
            \item {[x] Criar seis issues;}
            \item {[x] Reagir e comentar uma issue cadastrada por outra pessoa da equipe. O comentário deve conter informação útil para a possível solução da issue;}
            \item {[x] Revisar uma issue em status "closed"\ de outra pessoa;}
            \item {[ ] Registrar como comentário na issue revisada, os ajustes que precisam ser feitos.}
        \end{enumerate}
    
    \subsubsection{Controle de Versões}
        \begin{enumerate}
            \item {[ ] Criar um arquivo dentro do projeto;} 
            \item {[x] Alterar um arquivo dentro do projeto;}
            \item {[x] Revisar e aprovar um pull-request;}
            \item {[ ] Realizar uma alteração que resolva duas issues cadastradas por outro membro;}
            \item {[ ] Resolver, pelo menos, dois conflitos de merge.}
        \end{enumerate}
\subsection{Othon}

\subsection{Rafaela Santos S. Orrico}
\subsubsection{Gerenciamento de Mudanças}
    \begin{enumerate}
        \item {[x] Criar seis issues;}
        \item {[x] Reagir e comentar uma issue cadastrada por outra pessoa da equipe. O comentário deve conter informação útil para a possível solução da issue;}
        \item {[x] Revisar uma issue em status "closed"\ de outra pessoa;}
        \item {[ ] Registrar como comentário na issue revisada, os ajustes que precisam ser feitos.}
    \end{enumerate}
    
\subsubsection{Controle de Versões}
    \begin{enumerate}
        \item {[x] Criar um arquivo dentro do projeto;}
        \item {[x] Alterar um arquivo dentro do projeto;}
        \item {[ ] Revisar e aprovar um pull-request}
        \item {[ ] Realizar uma alteração que resolva duas issues cadastradas por outro membro;}
        \item {[ ] Resolver, pelo menos, dois conflitos de merge.}
    \end{enumerate}

\subsection{Tuanny}

\section{Evidências}
Esta seção terá as evidencias feita nesse trabalho, usando o git log.

\end{document}